\section{The action bound}
\label{sec:action-limit}

We now prove Theorem~\ref{thm:action} by combining the paths of
Proposition~\ref{prop:paths} with the field comparison.  This follows
the action identity and penalty argument in
\cite[Section~6]{ScalarCompanion}.  The path
equation produces a negative square in the action identity.  As
$\varepsilon$ tends to zero, this term forces any matrix values that
could violate the desired bound toward the constraint
$M_\delta(B)=AA^{\mathsf T}$.  At that constraint the field is positive
semidefinite, so Proposition~\ref{prop:field-comparison} applies.

\begin{proof}[Proof of Theorem~\ref{thm:action}]
Fix $\delta>0$ and, for each $0<\varepsilon\leq1$, choose
$C_\varepsilon$, $A_\varepsilon=C_\varepsilon U$, and
$B_\varepsilon$ as in Proposition~\ref{prop:paths}.  Temporarily
suppress the subscript $\varepsilon$.  Since $A\in H^1_0$ and
$B\in C^1$, the function
\[
    x\longmapsto J_\delta(B(x))+\tr(A(x)^{\mathsf T}B(x)A(x))
\]
is absolutely continuous.  The primitive identity
\eqref{eq:field-primitive} and the path equation give, almost
everywhere,
\begin{align*}
    \bigl(J_\delta(B)+\tr(A^{\mathsf T}BA)\bigr)'
       &=2\tr((A')^{\mathsf T}BA)
          -\tr\bigl((M_\delta(B)-AA^{\mathsf T})B'\bigr)\\
       &=2\tr((A')^{\mathsf T}BA)
          -\varepsilon^{-1}
              \|M_\delta(B)-AA^{\mathsf T}\|_{\HS}^2\\
       &\leq\|A'\|_{\HS}^2+\tr(B^2AA^{\mathsf T})
          -\varepsilon^{-1}
              \|M_\delta(B)-AA^{\mathsf T}\|_{\HS}^2.
\end{align*}
The last line is the Hilbert--Schmidt Cauchy--Schwarz inequality followed
by $2ab\leq a^2+b^2$.  Since $A$ vanishes at both endpoints,
integration and the definition of the action yield
\begin{align}
\label{eq:action-penalty}
    J_\delta(-K)-J_\delta(K)
       \leq{}&\mathcal E(A)
       +\int_\Omega\Bigl[
            \tr\bigl((AA^{\mathsf T})^q\bigr)
            -\tr\bigl((K^2-B^2)AA^{\mathsf T}\bigr)
            \notag\\[-2pt]
       &\hspace{40mm}
            -\varepsilon^{-1}
                \|M_\delta(B)-AA^{\mathsf T}\|_{\HS}^2
                         \Bigr]\,dx.
\end{align}

We next bound the integral uniformly as $\varepsilon\downarrow0$.
The bounds in \eqref{eq:path-bounds}, together with the boundedness of
$U$, place all pairs
\[
    (A_\varepsilon(x),B_\varepsilon(x)),
    \qquad x\in\overline\Omega,\quad0<\varepsilon\leq1,
\]
in one compact set $\mathcal K_\delta\subset
\R^{n\times d}\times\Sym_n$.  On this set define the continuous
functions
\begin{align*}
    F(A,B)&=\tr\bigl((AA^{\mathsf T})^q\bigr)
               -\tr\bigl((K^2-B^2)AA^{\mathsf T}\bigr),\\
    r_\delta(A,B)&=\|M_\delta(B)-AA^{\mathsf T}\|_{\HS}^2.
\end{align*}
If $r_\delta(A,B)=0$, then $M_\delta(B)=AA^{\mathsf T}\geq0$,
and Proposition~\ref{prop:field-comparison} gives $F(A,B)\leq E_\delta$.
This is the only point at which positivity of the field is needed;
the constructed paths themselves need not have this property.

Compactness now implies
\begin{equation}
\label{eq:penalty-limit}
    \limsup_{\varepsilon\downarrow0}
       \sup_{(A,B)\in\mathcal K_\delta}
          \bigl(F(A,B)-\varepsilon^{-1}r_\delta(A,B)\bigr)
       \leq E_\delta.
\end{equation}
Indeed, a failure would give $\eta>0$, a sequence
$\varepsilon_j\downarrow0$, and points $(A_j,B_j)\in\mathcal K_\delta$
at which the expression exceeds $E_\delta+\eta$.  Boundedness of $F$
then forces $r_\delta(A_j,B_j)=O(\varepsilon_j)$.  A convergent
subsequence has limit $(A_*,B_*)$ with $r_\delta(A_*,B_*)=0$ and,
by continuity, $F(A_*,B_*)\geq E_\delta+\eta$.  This contradicts
the preceding comparison.

Apply \eqref{eq:penalty-limit} in \eqref{eq:action-penalty}, bound
$\mathcal E(A_\varepsilon)$ by the supremum over lower triangular
$C$, and let $\varepsilon\downarrow0$.  We obtain
\begin{equation}
\label{eq:regularized-action-bound}
    J_\delta(-K)-J_\delta(K)
       \leq\sup_{C\in\mathcal L}\mathcal E(CU)
                       +|\Omega|E_\delta.
\end{equation}
All compactness estimates so far were for fixed $\delta$.  This is
sufficient: the supremum on the right is independent of $\delta$,
and the remaining limit concerns only the explicit endpoint values
and the error term.

By \eqref{eq:field-endpoints},
\[
    J_\delta(-K)-J_\delta(K)
       =\sum_{i=1}^n\int_{-k_i}^{k_i}
               \mu_\delta(k_i^2-s^2)\,ds.
\]
For $z\geq0$, the formula
$\mu_\delta(z)=(z+\delta)^\sigma-\delta^\sigma$ gives
\[
    0\leq\mu_\delta(z)\leq z^\sigma,
    \qquad \mu_\delta(z)\longrightarrow z^\sigma
            \quad\text{as }\delta\downarrow0.
\]
Dominated convergence therefore sends the endpoint difference to
\[
    \sum_{i=1}^n\int_{-k_i}^{k_i}(k_i^2-s^2)^\sigma\,ds.
\]
The error $E_\delta$ in \eqref{eq:field-error} tends to zero.  Letting
$\delta\downarrow0$ in \eqref{eq:regularized-action-bound} proves
the finite-interval action inequality.
\end{proof}
